\documentclass[../main.tex]{subfiles}
\begin{document}
\section{E50実験への適用}
\label{sec:methods}
4. MARQ実験への適用
30 Mcpsのビームへの適用

有効そうなフィルタ手法
\begin{itemize}
  \item SFTとT0とBFTだけでも散乱or not
  \item 60 Mcpsを真面目にトラッキング
\end{itemize}

ビーム下流のDCに散乱トラックがあるかないか

30 Mcps x 反応率 x 生成粒子数

スケーリング(T103の1 Mcps->MARQの30 Mcps(spill on avg))

BFTのトラッキングのスケーリング

タイミングカウンタのTOFのスケーリング

T0(ターゲット直前, もっと上流)

+α ビームのマイクロストラクチャの有無

+α シミュレーション


\subsection{T103データレートに対する負荷換算}
T103の入力データレートをspill on中546\,MB/s、ビーム取り出し時間を2\,s、休止時間を2.2\,sとすると、サイクル平均データレートは
\begin{equation}
  R_{\mathrm{avg}}=546\times\frac{2}{2+2.2}=260\,\mathrm{MB/s}
\end{equation}
となる\cite{kashima:interim2026}。
この平均入力レートを表\ref{tab:hlf-performance}の1プロセスあたりの処理性能で割ると、TOF、クラスター数、ローカルトラッキングについて、それぞれ0.44、0.59、2.11プロセス相当の負荷となる。
これらは平均入力と処理性能の比による換算値であり、整数の実装プロセス数やspill内のピーク負荷を直接表すものではない。

E50への適用では、ビームレートだけでなく、イベントウィンドウ内のヒット密度、検出器数、ヒットの組み合わせ、トリガー条件の変化を考慮する必要がある。
T103で得た性能をビームレートの比だけで線形外挿して、本番の必要計算資源とすることはできない。

\end{document}
